\section{Síntesis y ampliación}

El valor central de Meta-RL está en lograr una \textbf{adaptación rápida}:

\begin{enumerate}
    \item El RL multitarea aprende una política que cubre todas las tareas; meta-RL aprende la capacidad de adaptarse rápidamente a tareas nuevas
    \item Los métodos context-based realizan una inferencia implícita al codificar la experiencia histórica
    \item Los métodos gradient-based (MAML) logran un fine-tuning rápido al aprender una buena inicialización
    \item En teoría, meta-RL equivale a tomar decisiones óptimas en un POMDP con incertidumbre sobre la tarea
    \item El in-context learning de los LLM es una implementación a gran escala de meta-learning
\end{enumerate}

\subsection{Mapa de métodos de esta clase}

\begin{table}[H]
\centering
\begin{tabular}{p{0.18\textwidth} p{0.22\textwidth} p{0.22\textwidth} p{0.22\textwidth}}
\toprule
Enfoque & Mecanismo de adaptación & Ventajas & Dificultades \\
\midrule
Multi-task RL & task descriptor $z$ explícito & Estructura sencilla; admite compartir parámetros & Carece de un mecanismo de adaptación rápida ante tareas nuevas \\
Black-box Meta-RL & Contexto de experiencia + red de memoria & Gran capacidad expresiva; unificado de extremo a extremo & Difícil de entrenar y de explorar; la eficiencia de muestreo está limitada por el outer RL \\
Gradient-based Meta-RL & Aprender una inicialización que admita un fine-tuning rápido & Mecanismo de adaptación claro & Alto costo de la optimización de segundo orden; poca estabilidad en escenarios de RL \\
Structured exploration (PEARL, entre otros) & Muestreo posterior / recompensa intrínseca / modelado de la tarea & Facilita hacer explícita la estructura de la exploración & Puede introducir supuestos fuertes; subóptimo en algunos entornos \\
\bottomrule
\end{tabular}
\caption{Relación entre los métodos centrales de la Lecture 13}
\end{table}

\subsection{Lista de verificación para implementar Black-Box Meta-RL}

Al implementar en la práctica este tipo de algoritmos, lo que más se presta a errores no son las fórmulas, sino los «límites de la tarea» y los «límites de la memoria». Si el código de entrenamiento mezcla episode, task, hidden state y replay buffer, lo que se termina aprendiendo suele ser solo una recurrent policy común, y no una meta-policy capaz de adaptarse entre tareas.

\begin{table}[H]
\centering
\begin{tabular}{p{0.2\textwidth} p{0.28\textwidth} p{0.28\textwidth}}
\toprule
Módulo & Error frecuente & Práctica más segura \\
\midrule
Task sampling & Los límites entre tareas dentro de un mismo lote no están claros & Registrar explícitamente el task id y los límites de los episodes de cada tarea \\
Hidden state & Reutilizar por error el estado oculto entre tareas distintas & Conservar la memoria solo entre los episodes sucesivos de una misma tarea \\
Reward input & No introducir reward/history en la red como señal de adaptación & Usar explícitamente la secuencia $(s,a,r,s')$ como contexto \\
Meta-test protocol & Usar en la prueba las tareas de entrenamiento o filtrar la etiqueta de la tarea & Proporcionar solo pocos datos de interacción con la tarea nueva y evaluar estrictamente con held-out tasks fuera de la distribución \\
\bottomrule
\end{tabular}
\caption{Lista de verificación para implementar Black-box Meta-RL}
\end{table}

\begin{figure}[H]
\centering
\includegraphics[width=0.9\textwidth]{slides-images/slide-39.jpg}
\caption{Cierre de la clase: hoy se vieron los fundamentos de meta-RL; la próxima clase continúa con exploration y meta-exploration}
\end{figure}

\newpage
\subsection{Por qué la exploration se trata por separado}

Esta clase ya dio la respuesta: en el RL convencional, la exploración sirve para encontrar conductas de alta recompensa; en Meta-RL, la exploración además tiene la función de \textbf{identificar la tarea}. Por ello, una acción que «por ahora no parece rendir nada» puede ser justamente el paso de adaptación más importante. Precisamente por este doble papel, muchos métodos de black-box meta-RL, aunque conceptualmente unificados, se degradan de forma notable en entornos que exigen una exploración compleja.

El valor de discutir la exploration por separado es que nos recuerda no reducir Meta-RL a una «red de políticas con memoria». La parte realmente difícil es: \textbf{cómo lograr que el agent gaste primero un costo pequeño en averiguar qué tarea enfrenta y luego dedique el presupuesto restante a la estrategia de ejecución correcta.} Esa es también la razón de ser de toda la línea de investigación posterior, desde posterior sampling e intrinsic reward hasta task dynamics prediction.

\begin{knowledgebox}{Lecturas adicionales}
\begin{itemize}
    \item Finn et al., ``Model-Agnostic Meta-Learning for Fast Adaptation of Deep Networks (MAML),'' ICML 2017
    \item Duan et al., ``RL$^2$: Fast Reinforcement Learning via Slow Reinforcement Learning,'' ICLR 2017
    \item Rakelly et al., ``Efficient Off-Policy Meta-Reinforcement Learning via Probabilistic Context Variables (PEARL),'' ICML 2019
    \item Zintgraf et al., ``VariBAD: A Very Good Method for Bayes-Adaptive Deep RL,'' ICLR 2020
\end{itemize}
\end{knowledgebox}

\end{document}
